\subsection{DERIVATIONS OF AN ALGEBRA A INTO AN A-MODULE}

We suppose in this no. that Case (II) of no. 2 holds. Then there is a graded K-algebra A and a graded K-module E and also two K-linear mappings of degree 0

\[
\mathbf {E} \otimes_ {\mathbf {K}} \mathbf {A} \rightarrow \mathbf {E}, \quad \mathbf {A} \otimes_ {\mathbf {K}} \mathbf {E} \rightarrow \mathbf {E}
\]

denoted by

\[
x \otimes a \mapsto x. a \quad \text { and } \quad a \otimes x \mapsto a. x \quad \text { for   } a \in A \text {   and   } x \in E.
\]

PROPOSITION 3. Let \(d:A\to E\) be an \(\varepsilon\) -derivation of degree \(\delta\) . Then \(\operatorname{Ker}(d)\) is a graded subalgebra of A; if A admits a unit element, it belongs to \(\operatorname{Ker}(d)\) .

Clearly \(\operatorname{Ker}(d)\) is a graded sub-K-module of A; further, relation (8) of no. 2 shows that, if x and y are two homogeneous elements belonging to \(\operatorname{Ker}(d)\) , then \(d(xy) = 0\) and hence \(xy \in \operatorname{Ker}(d)\) . Finally, if A admits a unit element 1 (of degree 0, cf.~§3, no. 1), relation (8) of no. 2, where x and y are replaced by 1, gives \(d(1) = d(1) + d(1)\) and hence \(d(1) = 0\) .

COROLLARY. Let \(d_1\) and \(d_2\) be two \(\varepsilon\) -derivations from \(A\) to \(E\) of the same degree \(\delta\) . If \(d_1\) and \(d_2\) take the same values at each element of a generating system of the algebra \(A\) , then \(d_1 = d_2\) .

\(d_{1} - d_{2}\) is an \(\varepsilon\) -derivation of degree \(\delta\) , hence \(\operatorname{Ker}(d_1 - d_2)\) is a subalgebra of A which contains a generating system of A and hence is equal to A.

PROPOSITION 4. Let \(d:A\to E\) be an \(\varepsilon\) -derivation of degree \(\delta\) . Suppose that A has a unit element 1 and let \(x\) be a homogeneous element of A with an inverse \(x^{-1}\) in A. Then

\[
d (x ^ {- 1}) = - \varepsilon_ {\delta , \deg (x)} x ^ {- 1} ((d x) x ^ {- 1}) = - \varepsilon_ {\delta , \deg (x)} (x ^ {- 1} (d x)) x ^ {- 1}.\tag{19}
\]

We have \(d(xx^{-1}) = d(1) = 0\) (Proposition 3), whence

\[
(d x) x ^ {- 1} + \varepsilon_ {\delta , \deg (x)} x (d (x ^ {- 1})) = 0
\]

which proves the first formula of (19). On the other hand, \(x^{-1}\) is homogeneous of degree \(-\deg(x)\) and \(\varepsilon_{\delta,\deg(x)} = \varepsilon_{\delta,-\deg(x)}\) by formulae (1) of no. 1; writing \(d(x^{-1}x) = 0\) , the second formula of (19) is obtained similarly.

PROPOSITION 5. Suppose that A is an integral domain and let L be its field of fractions. Every derivation of A into a vector space E over L (considered as an A-module) can be extended uniquely to a derivation of L into E.

Let \(d\) be a derivation of \(\mathbf{A}\) into \(\mathbf{E}\) and \(\bar{d}\) a derivation of \(\mathbf{L}\) into \(\mathbf{E}\) extending \(d\) ; then, for \(u \in \mathbf{A}, v \in \mathbf{A}, v \neq 0\) , of necessity, by virtue of (19),

\[
\bar {d} (u / v) = v ^ {- 1} d u - u v ^ {- 2} d v\tag{20}
\]

which proves the uniqueness of \(\bar{d}\) . Conversely, we show that \(\bar{d}\) can be defined by formula (20); it must first be verified that if \(u / v = u' / v'\) the value of the right hand side of (20) does not change when \(u\) is replaced by \(u'\) and \(v\) by \(v'\) . Now, \(uv' = vu'\) , hence \(v'(du) + u(dv') = v(du') + u'(dv)\) and therefore \(v'(du - uv^{-1}dv) = v(du' - u'v'^{-1}dv')\) , since \(uv'v^{-1} = u'\) and \(u'v'^{-1}v = u\) . Thus a mapping \(\bar{d}: L \to E\) has been defined which extends \(d\) ; it is immediately verified that it is K-linear and a derivation.

PROPOSITION 6. Suppose that A is a unital associative graded K-algebra and E a graded (A, A)-bimodule. If \(d: \mathbf{A} \to \mathbf{E}\) is an \(\varepsilon\) -derivation of degree \(\delta\) , then, for every finite sequence \((x_i)_{1 \leqslant i \leqslant n}\) of homogeneous elements of A, of respective degrees \(\xi_i\) ( \(1 \leqslant i \leqslant n\) ),

\[
d \left(x _ {1} x _ {2} \dots x _ {n}\right) = \sum_ {i = 1} ^ {n} \varepsilon_ {\delta, \xi_ {1} + \dots + \xi_ {i - 1}} x _ {1} \dots x _ {i - 1} \left(d x _ {i}\right) x _ {i + 1} \dots x _ {n}.\tag{21}
\]

Formula (21) is trivial for \(n = 0\) and is proved by induction on \(n\) , taking account of (4) (no. 2).

COROLLARY. Suppose that A is a unital commutative associative algebra and E an A-module. If \(d: \mathbf{A} \to \mathbf{E}\) is a derivation, then, for every integer \(n \geqslant 0\) ,

\[
d (x ^ {n}) = n x ^ {n - 1} (d x) \quad \text { for all } x \in \mathbf {A}.\tag{22}
\]

It suffices to give A the trivial graduation and apply (21) with all the \(x_{i}\) equal to \(x\) .

We return to the general case of an \(\varepsilon\) -derivation \(d: \mathbf{A} \to \mathbf{E}\) of degree \(\delta\) . Let \(Z_{\varepsilon}\) be the set of \(a \in \mathbf{A}\) such that for every homogeneous component \(a_{\alpha}\) of \(a\) of degree \(\alpha\) , for every homogeneous \(x\) in \(\mathbf{E}\) ,

\[
x a _ {\alpha} = \varepsilon_ {\alpha , \deg (x)} a _ {\alpha} x.\tag{23}
\]

If A is a unital associative graded algebra and E a graded (A, A)-bimodule it follows immediately from this definition that \(Z_{\varepsilon}\) is a graded subalgebra of A containing the unit element.

PROPOSITION 7. Suppose that A is a unital associative graded algebra and E a graded (A, A)-bimodule. Let \(d: \mathbf{A} \to \mathbf{E}\) be an \(\varepsilon\) -derivation of degree \(\delta\) and \(a\) a homogeneous element of \(Z_{\varepsilon}\) of degree \(\alpha\) . Then the mapping \(x \mapsto a(dx)\) is an \(\varepsilon\) -derivation of degree \(\delta + \alpha\) .

We write \(d'(x) = a(dx)\) ; for \(x\) homogeneous of degree \(\xi\) in \(A\) and \(y \in A\) , by virtue of (23) and (1) (no. 1),

\[
\begin{array}{r l} d ^ {\prime} (x y) & = a ((d x) y) + \varepsilon_ {\delta , \xi} a (x (d y)) = (a (d x)) y + \varepsilon_ {\delta + \alpha , \xi} (x a) (d y) \\ & = (d ^ {\prime} x) y + \varepsilon_ {\delta + \alpha , \xi} x (d ^ {\prime} y). \end{array}
\]

Proposition 7 says that the K-module of \(\varepsilon\) -derivations of A into E is a graded \(Z_{\varepsilon}\) -module of type \(\Delta\) .

